% Pursuing Stacks (cotes 134-2 to 134-8): archive cote 134-4, site page 21 (PDF page 22), typescript page 109.
% typescript. Draft compared: ps3.tex:1025-1091.
% AI-transcribed by Claude Opus 5.5 (claude-opus-5-5, adaptive thinking, effort high), Message Batches API, batch ps2 (msgbatch_01NVwr1Yuo8qJf5rykeGBXPK), 2026-11-12; not yet checked by eye.
\scanpage{134-4/21}
\mspage{109}

In case when there exists a Lawvere element \xout{inxM} $L_M$ in M, \xout{thisxelemen}
for instance if M is a topos, this element is automatically endowed with
an idem\over{\unsure{b}}{p}otent composition law, coming from intersection of subobjects,
and in case M admits a strict initial object, $L_M$ is endowed with two
canonical sections which are indeed ``disjoint'', corresponding \over{\unsure{t}}{r}espective-%
ly to the ``full'' and the ``empty'' subobjects. T\over{\illegible{1}}{h}en (as already noticed
previously in the corollary of \deleted{page 60)} the comparison lemma) M admits a
homotopy interval iff $L_M$ is such an interval, namely is aspheric over
$e_M$ . But we can now add that in this case, there exists even a \uline{contract\unsure{ib|b}}%
\uline{le} homotopy interval, namely $L_M$ itself. \xout{\unsure{Thusxinxthisxcasex,xtherexxis}}
\margin{X} It is even a ``\uline{strict homotopy interval}'', namely one admitting a composi-%
tion $LxL \longrightarrow L$ having the properties above (where \hand{$\delta$}$_0$ and \hand{$\delta$}$_1$\hand{,} for the first
time\hand{,} play assymetric roles\hand{!}). Thus when a Lawvere element exists in M,
there is equivalence between the three \xout{conditionsx}\inserted[typed]{properties} we may expect from
a pair (M,W)\hand{,} with respect to homotopy intervals, namely :
\[ \hand{$($}\hand{$\exists$}\text{ homotopy intervals}\hand{$)\Leftarrow($}\hand{$\exists$}\text{ contractible hom. int}\hand{$)\Leftarrow($}\hand{$\exists$}\text{ strict hom. int.}\hand{$)$} . \]
It seems that in all cases I've in mind at present, when there exists a
homotopy interval, there exists even a strict one. T\over{\illegible{1}}{h}e only case besides
topoi (where it is s\over{\illegible{1}}{o}, because of the existence of a Lawvere element)
which I have looked \over{\illegible{1}}{up} so far, is the case of (Cat) and various full sub-%
\over{v}{c}ategories, all containing \hand{$\Delta$}$_1$ which is indeed a strict homotopy interval\hand{,}
as it represent\over{\unsure{g}}{s}s the \xout{functor} presheaf on (Cat)
\[ C \longmapsto \text{ set of all cribles in } C \ . \]
If we take the choice I = two-point category \over{\illegible{1}}{e}quivalent to the final one\hand{,}
this also is a strict homotopy interval, as it represents the functor
\[ C \longmapsto \text{ set of all full subcategories of } C \ , \]
hence again an intersection law. The first \inserted[typed]{homotopy} interval though is a lot more
important than the second, because the first one is ``characteristic''
namely sufficient for checking the homotopy relation between any two
maps in (Cat), whereas the second is'nt. A ``\uline{perfect}'' homotopy interval
would be one which is both \uline{strict} (hence contractible) and \uline{characteristic}.

2.4.

50. While writing the notes last time, and afterwards while pondering
a little more a\over{\illegible{1}}{b}out the matter, a few impressions came gradually into the
fore. One was about the interplay of four basic ``homotopy notions'' which
more or less mutually determine each other, namely the \uline{homotopy} \uline{relation}
between maps, the notion of \uline{homotopy} \uline{intervals}, the notion of homotopy
equivalences \over{a}{o}r \uline{homotopisms} (which has formal analogy to weak equivalences\hand{,}

\margin{21}
